\chapter{IgG4-Related Disease}
\index{IgG4-Related Disease}
\label{sec:IgG4-Related Disease}

\section{Overview}
IgG4-related disease is a fibroinflammatory condition characterized by tumefactive (mass-like) lesions, lymphoplasmacytic infiltration rich in IgG4-positive plasma cells, storiform scarring, and often elevated serum IgG4 levels, affecting virtually any organ with the most common sites being the pancreas (autoimmune pancreatitis type 1), salivary glands (sialadenitis, Mikulicz disease), orbits (dacryoadenitis, orbital pseudotumor), retroperitoneum (retroperitoneal scarring), thyroid (Riedel thyroiditis), kidneys (tubulointerstitial nephritis), lymph nodes, lungs, and biliary tree, with an incidence of 1--5 per 100,000, a male predominance (3:1), and typically presenting in the sixth decade.
\section{Causes}
This condition happens because of a complex interplay of genetic susceptibility, adaptive and innate immune activation, with a Th2-skewed immune response, T follicular helper cells, and overproduction of IL-4, IL-5, IL-13, and TGF-$\beta$, driving IgG4 class switching, eosinophilia, and scarring.     
\section{Risk Factors}

\begin{itemize}[noitemsep]
  \item Genetic predisposition and family history
  \item Advancing age
  \item Lifestyle factors (diet, exercise, smoking, alcohol)
  \item Environmental exposures
  \item Underlying medical conditions
  \item Medication use
\end{itemize}
\section{Signs and Symptoms}

\subsection{Common symptoms}
\begin{itemize}[noitemsep]
  \item Clinical features are organ-dependent but commonly include painless swelling of affected organs, obstructive symptoms (jaundice from pancreatic/biliary involvement, hydronephrosis from retroperitoneal scarring), and constitutional symptoms (fatigue, weight loss, low-grade fever). 
  \end{itemize}

\subsection{Emergency warning signs}
\begin{itemize}[noitemsep]
  \item Severe or worsening symptoms of igg4-related disease require immediate medical evaluation
\end{itemize}
\section{Progression and Stages}
The condition may progress through different stages of severity over time. Early stages often involve mild or subtle symptoms that may be overlooked. Moderate stages involve more noticeable symptoms affecting daily function. Advanced stages may involve significant impairment and complications.
\section{Complications}
\begin{itemize}[noitemsep]
  \item Disease progression leading to worsening symptoms
  \item Secondary infections or injuries
  \item Side effects of treatment
  \item Reduced quality of life
  \item Emotional and psychological impact
\end{itemize}
\section{Diagnosis}
To make the diagnosis, doctors look for histologic confirmation: characteristic histopathology (dense lymphoplasmacytic infiltrate with $>$10 IgG4+ plasma cells per high-power field, IgG4+/IgG+ ratio $>$40\%, storiform scarring, obliterative phlebitis). Comprehensive medical history and physical examination Laboratory tests to assess organ function and rule out other conditions Imaging studies to visualize affected structures Specialized testing depending on the affected system Consultation with relevant specialists Monitoring of disease progression over time

\begin{itemize}[noitemsep]
  \item Comprehensive medical history and physical examination
  \item Laboratory tests to assess organ function
  \item Imaging studies to visualize affected structures
  \item Specialized testing depending on the affected system
  \item Consultation with relevant specialists
\end{itemize}
\section{Prevention}
\begin{itemize}[noitemsep]
  \item Maintain a healthy lifestyle with balanced diet and exercise
  \item Avoid known risk factors
  \item Attend regular medical check-ups for early detection
  \item Follow recommended screening guidelines
\end{itemize}
\section{Treatment Options}
Treatment typically involves with corticosteroids (first-line, highly effective), with rituximab, azathioprine, mycophenolate, and methotrexate used for steroid-refractory or relapsing disease. Medications to manage symptoms and address underlying mechanisms Lifestyle modifications and self-care measures Physical therapy and rehabilitation Surgical intervention when indicated Regular monitoring and follow-up care Patient education and support services

\begin{itemize}[noitemsep]
  \item Medications to manage symptoms and address underlying mechanisms
  \item Lifestyle modifications and self-care measures
  \item Physical therapy and rehabilitation
  \item Surgical intervention when indicated
  \item Regular monitoring and follow-up care
\end{itemize}
\section{Living with It}
\begin{itemize}[noitemsep]
  \item Follow your treatment plan as prescribed
  \item Maintain regular follow-up with your healthcare provider
  \item Make necessary lifestyle adjustments
  \item Seek support from family, friends, or support groups
  \item Stay informed about your condition
\end{itemize}
\section{Outlook}
In most cases, the outlook is good but relapses are common (30--50\%) during steroid tapering. The outlook varies depending on the specific condition, severity at diagnosis, and response to treatment. Early intervention generally leads to better outcomes. Many conditions can be effectively managed with appropriate treatment. Regular follow-up care is important for monitoring and treatment adjustment.
